% chap1_introduccion.tex

\chapter{Introducción}   % capítulos numerados normales

\section{Contexto y motivación}
Los vehículos aéreos no tripulados (UAV, del inglés Unmanned Aerial Vehicle) han experimentado un crecimiento exponencial durante las últimas décadas, evolucionando desde aplicaciones militares, hasta convertirse en una herramienta de uso generalizado en sectores  civiles. comerciales y de investigación. En la actualidad, sus aplicaciones abarcan la agricultura de precisión, la inspección de infraestructuras, las operaciones de búsqueda y rescate y la cinematografía entre muchas otras\cite{UAV applications}.

El elemento central de cualquier operación con UAV es la estación de control en tierra (GCS, Ground Control Station), la interfaz software a través de la cual el operador monitoriza el estado del vehículo y controla el vuelo.  Históricamente, las GCS se han desarrollado como aplicaciones de escritorio fuertemente vinculadas a un sistema operativo concreto. Herramientas como Mission Planner\cite{missionplanner} y QGroundControl\cite{qgroundcontrol} representan el estándar de facto en código abierto, pero presentan una barrera de accesibilidad significativa. Requieren instalación local, no están preparadas para abrirse en web o móvil, y presentan una curva de aprendizaje pronunciada para operadores con poca experiencia.

A medida que las tecnologías web y móvil han crecido, ha surgido la oportunidad de construir interfaces de control ligeras y multiplataforma, que se ejecuten directamente en el navegador sin necesidad de instalación local. En contextos operativos, la posibilidad de conectarse a un dron desde cualquier dispositivo con navegador, o de presentar una interfaz sencilla de usar y entender, supone una mejora sustancial. 

Flutter\cite{flutter}, el framework de interfaz de usuario de código abierto desarrollado por Google, es un candidato adecuado para este tipo de aplicaciones. Permite apuntar a Android, iOS, web y escritorio, desde una única base de código, y su modelo de gestión de estados lo hace idóneo para interfaces que deben responder de forma continua a flujos de telemetría y reflejar cambios de estado del vehículo en tiempo real. A pesar de su creciente adopción en proyectos de IoT, existe una gran escasez de ejemplos documentados y reutilizables orientados específicamente al control de UAV. Los desarrolladores que deseen construir este tipo de interfaces tienen que crear soluciones a partir de ejemplos dispersos de dominios no relacionados, dificultando así los retos propios del control de vehículos en tiempo real.

Este proyecto nace para cubrir precisamente ese vacío mediante dos contribuciones complementarias, desarrolladas en el marco del DEE (Drone Engineering Ecosystem)\cite{dronsEETAC}. La primera es una colección de aplicaciones Flutter de demostración, independientes entre sí, cada una explorando un paradigma de control de dron específico: control RC mediante un doble joystick virtual, reconocimiento de comandos por voz, control inercial mediante la IMU del propio dispositivo, localización GPS del operador, streaming de vídeo con detección de objetos, planificación de misiones autónomas y registro y análisis de vuelos. Cada demostración está documentada de forma autónoma y es accesible públicamente como referencia para futuros desarrolladores del ecosistema.

La segunda contribución es EZDrone, una estación de control en tierra basada en Flutter Web que integra todos estos paradigmas en una arquitectura coherente. EZDrone permite controlar dos plataformas de vuelo presentes en el DEE. El dron principal, un chasis HexSon con controlador de vuelo Pixhawk\cite{pixhawk} ejecutando firmware ArduPilot\cite{ardupilot}, utilizado para operaciones en exteriores, y el dron DJI Tello\cite{tello}, usado para interiores, especialmente en demostraciones a visitantes. 

En conjunto, este proyecto amplía el bloque de aplicaciones Flutter del DEE con recursos reutilizables y una estación tierra web accesible desde cualquier dispositivo y cualquier navegador, orientada tanto al ecosistema abierto ArduPilot como al uso docente y demostrativo de la EETAC.



\section{Estado del arte}
\subsection{Estaciones de control en tierra de código abierto}
El ecosistema de UAV de código abierto está estructurado en torno a dos familias de autopilot, ArduPilot y PX4\cite{px4},  que emplean MAVLink\cite{mavlink} como protocolo de comunicación principal. Sobre este sustrato se han construido las principales GCS disponibles hoy en día.

Mission Planner es la  GCS de referencia para ArduPilot, desarrollada en C\#. Cubre el ciclo operativo completo: configuración de parámetros, calibración de sensores, planificación de misiones, monitorización de telemetría en tiempo real, control completo del UAV y análisis de logs post-vuelo. Es la herramienta elegida por la mayoría de los usuarios de ArduPilot, sin embargo, está limitada exclusivamente a Windows, y no contempla ninguna forma de despliegue web o móvil.

QGroundControl amplía significativamente el soporte de plataformas, es compatible con Windows, macOS, Linux, Android e iOS. Ofrece una interfaz más moderna y accesible. A pesar de sus versiones móviles, su funcionalidad está muy reducida respecto a la versión de escritorio, y no existe ninguna versión web.

La documentación oficial de ArduPilot identifica al menos diez GCS, pero ninguna de las opciones ofrece actualmente despliegue web\cite{ardupilot_gcs}. Esta ausencia es la brecha directa que EZDrone pretende cubrir.   


\subsection{Aplicaciones móviles para el control de drones}
En el sector comercial, fabricantes como DJI y Skydio han desarrollado aplicaciones propias que ofrecen una experiencia de usuario muy cuidada e integración completa con su hardware. Aplicaciones como DJI Fly\cite{djifly} presentan una interfaz simple y sencilla que cuenta con funciones de control de vuelo, transmisión de vídeo en tiempo real, modos de vuelo automatizados (como seguimiento de personas, vuelos orbitales, panorámicas) y edición de contenido directamente en el dispositivo. Skydio\cite{skydio}, por su parte, pone el foco en la autonomía de vuelo y la evasión de obstáculos, ofreciendo interfaces pensadas para que el operador pueda realizar vuelos complejos sin requerir de mucha experiencia y con la mínima intervención manual. Estas soluciones están muy pulidas desde el punto de vista de la experiencia de uso, es decir, tienen interfaces muy intuitivas y guiadas, y ofrecen conectividad fiable y funcional con los propios dispositivos del fabricante.

Sin embargo, este nivel de integración tiene un coste importante, es código cerrado, diseñado exclusivamente para su propio hardware. No existe ninguna interfaz de programación pública para adaptarlas a otros autopilots, protocolos de comunicación o plataformas de terceros. Esto las hace irrelevantes en entornos académicos como es el caso del DEE, donde la flexibilidad y reutilización del código son requisitos fundamentales.

Este proyecto, con EZDrone y las demostraciones, adopta un enfoque totalmente distinto. Al estar construido sobre Flutter Web y formar parte del DEE, es independiente del hardware (quitando el caso del Tello en EZDrone). Puede conectarse a cualquier plataforma que ejecute firmware ArduPilot, o extenderse para soportar otros autopilots. Las demostraciones son proyectos Flutter autónomos y documentados, pensados para ser estudiados, modificados e integrados para futuros desarrolladores. la cotrapartida es que el nivel de pulido visual es mucho menor que en una solución comercial dedicada. EZDrone está orientada más al control funcional, la docencia e investigación, no tanto a la experiencia de usuario del gran público. Esta diferencia de objetivos es la que justifica y delimita el enfoque adoptado en este trabajo.


\subsection{Flutter en IoT y robótica}
Flutter ha consolidado su presencia en el desarrollo IoT desde su lanzamiento en 2018\cite{flutter}, principalmente en cuadros de mando de sensores y aplicaciones de domótica que se comunican mediante MQTT. La combinación Flutter + MQTT\cite{mqtt_dart} cuenta con documentación razonable para casos de uso sencillos: lectura de valores de sensores, visualización de estado y envío de comandos discretos.

Su uso para el control de vehículos en tiempo real introduce retos cualitativamente distintos que están mucho menos documentados. El control de un UAV exige generar salidas de forma continua (como los valores de joystick enviados), gestionar múltiples flujos asíncronos simultáneamente (telemetría, vídeo, comandos de estado) y mantener la coherencia de la interfaz ante cambios de estado rápidos del vehículo. Trabajos recientes demuestran que arquitecturas similares (Flutter como frontend, MQTT para señalización y WebRTC para flujos de alta frecuencia) son técnicamente viables\cite{sun2025}. 

\section{TFG y TFM relacionados}
El control de drones mediante interfaces móviles y web, así como el desarrollo de aplicaciones Flutter, han sido objeto de Trabajos de Fin de Grado y Trabajos de Fin de Máster en los últimos años. Estos trabajos han explorado aspectos como la implementación de frontends para el control de UAV 

Este proyecto se enmarca en esa línea de trabajos, pero con un foco específico y diferenciado, la creación de una colección de demostraciones Flutter reutilizables que cubran los principales paradigmas de control de dron, y el desarrollo de EZDrone como estación de control web integrada en el DEE. 


\section{Estructura del documento}
El Capítulo 2 presenta el Drone Engineering Ecosystem (DEE), la plataforma académica desarrollada en la EETAC dentro de la cual se encuentra este proyecto. Se describe su arquitectura general y los tipos de módulos que lo componen, los drones utilizados, y se explica la contribución de este trabajo al bloque de aplicaciones Flutter del ecosistema.

El Capítulo 3 agrupa el marco tecnológico del proyecto. Presenta el framework Flutter y el lenguaje de programación Dart, los modos de comunicación del DEE, las interfaces de programación de aplicaciones (API), y los protocolos de transporte y de aplicación usados (TCP/UDP, HTTP/HTTPS, MQTT y WebRTC). El capítulo concluye con los repositorios y herramientas de desarrollo utilizados a lo largo del proyecto,

El Capítulo 4 recoge los objetivos del proyecto y el plan de trabajo. Se detalla el objetivo general, los objetivos específicos y una descripción ordenada de las tareas principales a realizar.

El Capítulo 5 presenta la arquitectura software base capaz del control mínimo del dron. 

El capítulo 6 documenta cada una de las demostraciones de paradigmas de control desarrolladas de forma independiente junto a su implementación en EZDrone y una demostración de su funcionamiento. 

El Capítulo 7 describe la evolución de la arquitectura de comunicación, desde el diseño inicial (Capítulo 5) hasta el diseño híbrido definitivo. 

El Capítulo 8 presenta EZDrone en su conjunto, detallando la arquitectura, los detalles de implementación y los resultados de las pruebas. 

El Capítulo 9 recoge el estudio de la sostenibilidad. 

El Capítulo 10 cierra el documento con las conclusiones y las líneas de trabajo futuro.

